% GENERATED FILE. Edit canonical lessons, then run npm run generate:books.
% canonical-source-hash: fnv1a64:fce071a0c514c6d6
% canonical-lessons: SA-C391-paramah, SA-C391-nikrstah, SA-C391-anurupah, SA-C391-upayuktah, SA-C391-dusayati

\chapter{Better and Worse}
\label{ch:sa-better-and-worse}
\hlchaptermodality{391}

\begin{chapteropening}
\textbf{By the end of this chapter:} \emph{I can say whether things will get better or worse.}

Complete the last lesson of chapter 391: I can say whether things will get better or worse.
\end{chapteropening}

\section[\texorpdfstring{\sk{परमः} (paramaḥ)}{परमः (paramaḥ)}]{\sk{परमः} (paramaḥ) — supreme; highest}

\label{lesson:SA-C391-paramah}



\begin{quote}
Before the new one: say the Sanskrit for worthwhile, then the Sanskrit for experiences.
\end{quote}

\subsection*{You'll want to know: \sk{परमः}}
\textbf{\sk{परमः}} — \emph{paramaḥ} — \textquotedblleft{}supreme; highest\textquotedblright{}.

\begin{quote}
\sk{यदि} \sk{मनः} \sk{शान्तम्}, \sk{तर्हि} \sk{परमं} \sk{सुखं} \sk{भवति।}
\end{quote}

\begin{quote}
\emph{yadi manaḥ śāntam, tarhi paramaṃ sukhaṃ bhavati.}
\end{quote}

\begin{quote}
If the mind is calm, there is the highest happiness.
\end{quote}

\begin{cousinweb}
\textbf{\sk{परमः}} (\emph{paramaḥ}) is built on \emph{para}, \textquotedblleft{}far, beyond\textquotedblright{}, with the superlative ending \emph{-ma}; MW gives \emph{parama} \textquotedblleft{}chief, highest, primary, best, most excellent\textquotedblright{} (MW).
\end{cousinweb}

\begin{grammarlens}[title={What you've built}]
One more word for saying what happens if.
\end{grammarlens}

\subsection*{Your turn}
\begin{itemize}
\raggedright
  \item \emph{Say it:} \emph{paramaḥ}
  \item \emph{Say it:} \emph{paramaḥ}, once more
  \item \emph{Say it:} \emph{yadi manaḥ śāntam, tarhi paramaṃ sukhaṃ bhavati.} — If the mind is calm, there is the highest happiness.
  \item \emph{Recall:} say the Sanskrit for worthwhile, then the Sanskrit for experiences, then say \emph{paramaḥ} again
\end{itemize}

\begin{tcolorbox}[breakable,colback=teal!4,colframe=teal!35!black,title={Before you move on}]
What does \sk{परमः} mean? (\textquotedblleft{}Supreme; highest\textquotedblright{}.) Say it once more.
\end{tcolorbox}
\section[\texorpdfstring{\sk{निकृष्टः} (nikṛṣṭaḥ)}{निकृष्टः (nikṛṣṭaḥ)}]{\sk{निकृष्टः} (nikṛṣṭaḥ) — inferior; low}

\label{lesson:SA-C391-nikrstah}



\begin{quote}
Before the new one: say the Sanskrit for experiences, then the Sanskrit for supreme.
\end{quote}

\subsection*{You'll want to know: \sk{निकृष्टः}}
\textbf{\sk{निकृष्टः}} — \emph{nikṛṣṭaḥ} — \textquotedblleft{}inferior; low\textquotedblright{}.

\begin{quote}
\sk{यदि} \sk{वस्तु} \sk{निकृष्टम्}, \sk{तर्हि} \sk{तत्} \sk{न} \sk{क्रीणातु।}
\end{quote}

\begin{quote}
\emph{yadi vastu nikṛṣṭam, tarhi tat na krīṇātu.}
\end{quote}

\begin{quote}
If an article is inferior, do not buy it.
\end{quote}

\begin{cousinweb}
\textbf{\sk{निकृष्टः}} (\emph{nikṛṣṭaḥ}) is built on \emph{ni}, \textquotedblleft{}down\textquotedblright{}, and the root √kṛṣ, \textquotedblleft{}to draw\textquotedblright{}; MW gives \emph{ni-kṛṣṭa} \textquotedblleft{}debased, vile, low, despised\textquotedblright{} (MW).
\end{cousinweb}

\begin{grammarlens}[title={What you've built}]
One more word for saying what happens if.
\end{grammarlens}

\subsection*{Your turn}
\begin{itemize}
\raggedright
  \item \emph{Say it:} \emph{nikṛṣṭaḥ}
  \item \emph{Say it:} \emph{nikṛṣṭaḥ}, once more
  \item \emph{Say it:} \emph{yadi vastu nikṛṣṭam, tarhi tat na krīṇātu.} — If an article is inferior, do not buy it.
  \item \emph{Recall:} say the Sanskrit for experiences, then the Sanskrit for supreme, then say \emph{nikṛṣṭaḥ} again
\end{itemize}

\begin{tcolorbox}[breakable,colback=teal!4,colframe=teal!35!black,title={Before you move on}]
What does \sk{निकृष्टः} mean? (\textquotedblleft{}Inferior; low\textquotedblright{}.) Say it once more.
\end{tcolorbox}
\section[\texorpdfstring{\sk{अनुरूपः} (anurūpaḥ)}{अनुरूपः (anurūpaḥ)}]{\sk{अनुरूपः} (anurūpaḥ) — in keeping with; suitable}

\label{lesson:SA-C391-anurupah}



\begin{quote}
Before the new one: say the Sanskrit for supreme, then the Sanskrit for inferior.
\end{quote}

\subsection*{You'll want to know: \sk{अनुरूपः}}
\textbf{\sk{अनुरूपः}} — \emph{anurūpaḥ} — \textquotedblleft{}in keeping with; suitable\textquotedblright{}.

\begin{quote}
\sk{यदि} \sk{वेतनं} \sk{कार्यस्य} \sk{अनुरूपम्}, \sk{तर्हि} \sk{कर्मकरः} \sk{तुष्यति।}
\end{quote}

\begin{quote}
\emph{yadi vetanaṃ kāryasya anurūpam, tarhi karmakaraḥ tuṣyati.}
\end{quote}

\begin{quote}
If the wage fits the work, the worker is content.
\end{quote}

\begin{cousinweb}
\textbf{\sk{अनुरूपः}} (\emph{anurūpaḥ}) is built from \emph{anu}, \textquotedblleft{}after\textquotedblright{}, and \emph{rūpa}, \textquotedblleft{}form\textquotedblright{}; MW gives \emph{anu-rūpa} \textquotedblleft{}following the form, conformable, corresponding, fit, suitable\textquotedblright{} (MW).
\end{cousinweb}

\begin{grammarlens}[title={What you've built}]
One more word for saying what happens if.
\end{grammarlens}

\subsection*{Your turn}
\begin{itemize}
\raggedright
  \item \emph{Say it:} \emph{anurūpaḥ}
  \item \emph{Say it:} \emph{anurūpaḥ}, once more
  \item \emph{Say it:} \emph{yadi vetanaṃ kāryasya anurūpam, tarhi karmakaraḥ tuṣyati.} — If the wage fits the work, the worker is content.
  \item \emph{Recall:} say the Sanskrit for supreme, then the Sanskrit for inferior, then say \emph{anurūpaḥ} again
\end{itemize}

\begin{tcolorbox}[breakable,colback=teal!4,colframe=teal!35!black,title={Before you move on}]
What does \sk{अनुरूपः} mean? (\textquotedblleft{}In keeping with; suitable\textquotedblright{}.) Say it once more.
\end{tcolorbox}
\section[\texorpdfstring{\sk{उपयुक्तः} (upayuktaḥ)}{उपयुक्तः (upayuktaḥ)}]{\sk{उपयुक्तः} (upayuktaḥ) — suitable; useful}

\label{lesson:SA-C391-upayuktah}



\begin{quote}
Before the new one: say the Sanskrit for inferior, then the Sanskrit for in keeping with.
\end{quote}

\subsection*{You'll want to know: \sk{उपयुक्तः}}
\textbf{\sk{उपयुक्तः}} — \emph{upayuktaḥ} — \textquotedblleft{}suitable; useful\textquotedblright{}.

\begin{quote}
\sk{यदि} \sk{वस्त्रम्} \sk{उपयुक्तम्}, \sk{तर्हि} \sk{उत्सवे} \sk{धारयतु।}
\end{quote}

\begin{quote}
\emph{yadi vastram upayuktam, tarhi utsave dhārayatu.}
\end{quote}

\begin{quote}
If the garment is suitable, wear it at the festival.
\end{quote}

\begin{cousinweb}
\textbf{\sk{उपयुक्तः}} (\emph{upayuktaḥ}) is built on \emph{upa}, \textquotedblleft{}towards\textquotedblright{}, and the root √yuj, \textquotedblleft{}to join\textquotedblright{}; MW gives \emph{upa-yukta} \textquotedblleft{}suitable, fit, appropriate, useful\textquotedblright{} (MW).
\end{cousinweb}

\begin{grammarlens}[title={What you've built}]
One more word for saying what happens if.
\end{grammarlens}

\subsection*{Your turn}
\begin{itemize}
\raggedright
  \item \emph{Say it:} \emph{upayuktaḥ}
  \item \emph{Say it:} \emph{upayuktaḥ}, once more
  \item \emph{Say it:} \emph{yadi vastram upayuktam, tarhi utsave dhārayatu.} — If the garment is suitable, wear it at the festival.
  \item \emph{Recall:} say the Sanskrit for inferior, then the Sanskrit for in keeping with, then say \emph{upayuktaḥ} again
\end{itemize}

\begin{tcolorbox}[breakable,colback=teal!4,colframe=teal!35!black,title={Before you move on}]
What does \sk{उपयुक्तः} mean? (\textquotedblleft{}Suitable; useful\textquotedblright{}.) Say it once more.
\end{tcolorbox}
\section[\texorpdfstring{\sk{दूषयति} (dūṣayati)}{दूषयति (dūṣayati)}]{\sk{दूषयति} (dūṣayati) — spoils; pollutes}

\label{lesson:SA-C391-dusayati}



\begin{quote}
Before the new one: say the Sanskrit for in keeping with, then the Sanskrit for suitable.
\end{quote}

\subsection*{You'll want to know: \sk{दूषयति}}
\textbf{\sk{दूषयति}} — \emph{dūṣayati} — \textquotedblleft{}spoils; pollutes\textquotedblright{}.

\begin{quote}
\sk{यदि} \sk{यन्त्राणि} \sk{जलं} \sk{दूषयन्ति}, \sk{तर्हि} \sk{मत्स्याः} \sk{म्रियन्ते।}
\end{quote}

\begin{quote}
\emph{yadi yantrāṇi jalaṃ dūṣayanti, tarhi matsyāḥ mriyante.}
\end{quote}

\begin{quote}
If machines pollute the water, the fish die.
\end{quote}

Past (it …): \sk{धूमः} \sk{वायुं} \sk{दूषितवान्}, \emph{dhūmaḥ vāyuṃ dūṣitavān}.

\begin{cousinweb}
\textbf{\sk{दूषयति}} (\emph{dūṣayati}) comes from the root √duṣ, \textquotedblleft{}to become bad\textquotedblright{}, as a verb; MW gives \emph{dūṣaya} \textquotedblleft{}to corrupt, to spoil, to contaminate\textquotedblright{} and \emph{dūṣita} \textquotedblleft{}spoiled, corrupted\textquotedblright{} (MW).
\end{cousinweb}

\begin{grammarlens}[title={What you've built}]
One more word for saying what happens if.
\end{grammarlens}

\subsection*{Your turn}
\begin{itemize}
\raggedright
  \item \emph{Say it:} \emph{dūṣayati}
  \item \emph{Say it:} \emph{dūṣayati}, once more
  \item \emph{Say it:} \emph{yadi yantrāṇi jalaṃ dūṣayanti, tarhi matsyāḥ mriyante.} — If machines pollute the water, the fish die.
  \item \emph{Recall:} say the Sanskrit for in keeping with, then the Sanskrit for suitable, then say \emph{dūṣayati} again
\end{itemize}

\begin{tcolorbox}[breakable,colback=teal!4,colframe=teal!35!black,title={Before you move on}]
What does \sk{दूषयति} mean? (\textquotedblleft{}Spoils; pollutes\textquotedblright{}.) Say it once more.
\end{tcolorbox}
